跨数据集语义分割的得分接近零，本身不足以区分影像域偏移与源标注未被分配到待评估输出类别这两种情况。本文结合原始掩膜覆盖审计与固定预算的源数据混合实验考察这一问题，并明确采用七叶类别的操作性映射；其中，OpenEarthMap 的 rangeland 被用作植被代理类别，其原生定义也包含灌木。在主要的 LoveDA--OpenEarthMap 案例中，加入所考察的 FLAIR 草本源数据混合，使 rangeland 的 IoU 提高0.3118，但使六个 LoveDA 支持类别上的 mIoU 降低0.1011。对同一目标数据上的历史原生 ResNet-34 U-Net 记录进行再分析，得到相同方向：rangeland 恢复量为0.2828，支持类别 mIoU 的变化为-0.1304。反向配置的源混合在恢复耕地的同时也产生了支持类别性能损失；Semantic Drone 则提供了主要 CLIP 检查点在另一目标上的评估。这些结果通过特定的经验审计设计补充了已有标签空间与负迁移研究，没有提出新的分割算法或恢复指标。源混合同时改变影像和标签，且历史检查点与排程的来源证据仍不完整。因此，现有证据支持有边界的描述性对比，不能单独识别标签的因果作用，也不能证明已提出成功的缓解方法。正例源标签覆盖、缺失类别恢复与源支持类别保持应当共同报告。
